% !TEX root = ../main.tex
\begin{abstract}
Continuous diffusion offers a way to generate an entire reasoning trace through parallel refinement, but its effectiveness depends on the representation being denoised and on how the problem is supplied to the model. We develop a reasoning-oriented recipe based on Embedded Language Flows (ELF). We compress activations from three layers of a frozen Qwen teacher using projectors trained to preserve downstream token predictions, with covariance whitening to control latent geometry. A compact prompt Transformer first learns to imitate teacher prompt representations and is then trained jointly with a mature denoiser. Generation follows continuous layer-specific clocks and decodes to tokens only at the end. Before reward-based fine-tuning, our base and large denoisers achieve mean single-sample accuracies of 39.17\% and 58.67\% on GSM8K, and the large model achieves 18.85\% on MATH500. DiffusionNFT increases these to 42.16\%, 63.74\%, and 23.60\%, respectively. Eight samples yield 81.96\% and 52.60\% oracle accuracy for the large model on the two benchmarks. Six-epoch controls favor learned projectors, multilayer representations, and teacher prompt conditioning over joint learning from initialization. Immediate prompt substitution incurs a substantial loss that subsequent joint training largely recovers. Reward-based fine-tuning improves single-sample correctness, while its effects on coverage and answer diversity depend on the dataset. Finally, we show that the preferred allocation of inference compute changes when success is measured by majority voting rather than oracle pass@$k$.
\end{abstract}
